% GENERATED FILE. Edit canonical lessons, then run npm run generate:books.
% canonical-source-hash: fnv1a64:bbe48465a90cbba7
% canonical-lessons: TE-C341-mukhyamsalu, TE-C341-sampadakudu, TE-C341-sampadakiyam, TE-C341-pathakudu, TE-C341-dhruvikarincu, TE-R341-review-334-341

\chapter{Today's Headlines}
\label{ch:te-today-s-headlines}
\hlchaptermodality{341}

\begin{chapteropening}
\textbf{By the end of this chapter:} \emph{I can say the main points, an editor, an editorial, a reader and to confirm.}

Complete the last lesson of chapter 341: I can say the main points, an editor, an editorial, a reader and to confirm.
\end{chapteropening}

\section[\texorpdfstring{\te{ముఖ్యాంశాలు} (mukhyāṁśālu)}{ముఖ్యాంశాలు (mukhyāṁśālu)}]{\te{ముఖ్యాంశాలు} (mukhyāṁśālu) — the main points; highlights}

\label{lesson:TE-C341-mukhyamsalu}



\begin{quote}
Before the new one: say the Telugu for in the end, then the Telugu for to wake up.
\end{quote}

\subsection*{You'll want to know: \te{ముఖ్యాంశాలు}}
\textbf{\te{ముఖ్యాంశాలు}} — \emph{mukhyāṁśālu} — \textquotedblleft{}the main points; highlights\textquotedblright{}.

\begin{quote}
\te{రేడియోలో} \te{వార్తల} \te{ముఖ్యాంశాలు} \te{చెప్పారు}.
\end{quote}

\begin{quote}
\emph{rēḍiyōlō vārtala mukhyāṁśālu ceppāru.}
\end{quote}

\begin{quote}
The radio gave the main points of the news.
\end{quote}

\te{ముఖ్య}, chief, and \te{అంశాలు}, points: also the headlines read at the start of a bulletin. On election night the \te{ముఖ్యాంశాలు} are all about the \te{ఓటు}.

\begin{cousinweb}
\textbf{\te{ముఖ్యాంశాలు}} (\emph{mukhyāṁśālu}) is a compound borrowed from Sanskrit: \emph{mukhya}, \textquotedblleft{}chief\textquotedblright{} (CDIAL 10174), and \emph{aṁśa}, \textquotedblleft{}a share, a part\textquotedblright{} (CDIAL 2).
\end{cousinweb}

\begin{grammarlens}[title={What you've built}]
One more word for telling what happened.
\end{grammarlens}

\subsection*{Your turn}
\begin{itemize}
\raggedright
  \item \emph{Say it:} \emph{mukhyāṁśālu}
  \item \emph{Say it:} \emph{mukhyāṁśālu}, once more
  \item \emph{Say it:} \emph{rēḍiyōlō vārtala mukhyāṁśālu ceppāru.} — The radio gave the main points of the news.
  \item \emph{Recall:} say the Telugu for in the end, then the Telugu for to wake up, then say \emph{mukhyāṁśālu} again
\end{itemize}

\begin{tcolorbox}[breakable,colback=teal!4,colframe=teal!35!black,title={Before you move on}]
What does \te{ముఖ్యాంశాలు} mean? (\textquotedblleft{}The main points; highlights\textquotedblright{}.) Say it once more.
\end{tcolorbox}
\section[\texorpdfstring{\te{సంపాదకుడు} (sampādakuḍu)}{సంపాదకుడు (sampādakuḍu)}]{\te{సంపాదకుడు} (sampādakuḍu) — an editor}

\label{lesson:TE-C341-sampadakudu}



\begin{quote}
Before the new one: say the Telugu for to wake up, then the Telugu for the main points.
\end{quote}

\subsection*{You'll want to know: \te{సంపాదకుడు}}
\textbf{\te{సంపాదకుడు}} — \emph{sampādakuḍu} — \textquotedblleft{}an editor\textquotedblright{}.

\begin{quote}
\te{సంపాదకుడు} \te{ఆ} \te{శీర్షిక} \te{మార్చాడు}.
\end{quote}

\begin{quote}
\emph{sampādakuḍu ā śīrṣika mārcāḍu.}
\end{quote}

\begin{quote}
The editor changed that headline.
\end{quote}

Plural: \te{సంపాదకులు}, \emph{sampādakulu}.

The person who puts a paper together each day.

\begin{cousinweb}
\textbf{\te{సంపాదకుడు}} (\emph{sampādakuḍu}) is built on Sanskrit \emph{sampādaka}, \textquotedblleft{}producing, bringing about\textquotedblright{} (Monier-Williams), with the Telugu ending \emph{-uḍu}: the one who brings the paper about.
\end{cousinweb}

\begin{grammarlens}[title={What you've built}]
One more word for telling what happened.
\end{grammarlens}

\subsection*{Your turn}
\begin{itemize}
\raggedright
  \item \emph{Say it:} \emph{sampādakuḍu}
  \item \emph{Say it:} \emph{sampādakuḍu}, once more
  \item \emph{Say it:} \emph{sampādakuḍu ā śīrṣika mārcāḍu.} — The editor changed that headline.
  \item \emph{Recall:} say the Telugu for to wake up, then the Telugu for the main points, then say \emph{sampādakuḍu} again
\end{itemize}

\begin{tcolorbox}[breakable,colback=teal!4,colframe=teal!35!black,title={Before you move on}]
What does \te{సంపాదకుడు} mean? (\textquotedblleft{}An editor\textquotedblright{}.) Say it once more.
\end{tcolorbox}
\section[\texorpdfstring{\te{సంపాదకీయం} (sampādakīyaṁ)}{సంపాదకీయం (sampādakīyaṁ)}]{\te{సంపాదకీయం} (sampādakīyaṁ) — an editorial}

\label{lesson:TE-C341-sampadakiyam}



\begin{quote}
Before the new one: say the Telugu for the main points, then the Telugu for an editor.
\end{quote}

\subsection*{You'll want to know: \te{సంపాదకీయం}}
\textbf{\te{సంపాదకీయం}} — \emph{sampādakīyaṁ} — \textquotedblleft{}an editorial\textquotedblright{}.

\begin{quote}
\te{పత్రిక} \te{సంపాదకీయంలో} \te{ఆ} \te{సంఘటన} \te{గురించి} \te{రాసింది}.
\end{quote}

\begin{quote}
\emph{patrika sampādakīyaṁlō ā saṅghaṭana guriñci rāsindi.}
\end{quote}

\begin{quote}
The paper wrote about the incident in its editorial.
\end{quote}

The editor's own piece of writing: built on \te{సంపాదకుడు}.

\begin{cousinweb}
\textbf{\te{సంపాదకీయం}} (\emph{sampādakīyaṁ}) is a modern word built on Sanskrit \emph{sampādaka}, \textquotedblleft{}producing, bringing about\textquotedblright{} (Monier-Williams), with the Sanskrit ending \emph{-īya}, \textquotedblleft{}belonging to\textquotedblright{}.
\end{cousinweb}

\begin{grammarlens}[title={What you've built}]
One more word for telling what happened.
\end{grammarlens}

\subsection*{Your turn}
\begin{itemize}
\raggedright
  \item \emph{Say it:} \emph{sampādakīyaṁ}
  \item \emph{Say it:} \emph{sampādakīyaṁ}, once more
  \item \emph{Say it:} \emph{patrika sampādakīyaṁlō ā saṅghaṭana guriñci rāsindi.} — The paper wrote about the incident in its editorial.
  \item \emph{Recall:} say the Telugu for the main points, then the Telugu for an editor, then say \emph{sampādakīyaṁ} again
\end{itemize}

\begin{tcolorbox}[breakable,colback=teal!4,colframe=teal!35!black,title={Before you move on}]
What does \te{సంపాదకీయం} mean? (\textquotedblleft{}An editorial\textquotedblright{}.) Say it once more.
\end{tcolorbox}
\section[\texorpdfstring{\te{పాఠకుడు} (pāṭhakuḍu)}{పాఠకుడు (pāṭhakuḍu)}]{\te{పాఠకుడు} (pāṭhakuḍu) — a reader}

\label{lesson:TE-C341-pathakudu}



\begin{quote}
Before the new one: say the Telugu for an editor, then the Telugu for an editorial.
\end{quote}

\subsection*{You'll want to know: \te{పాఠకుడు}}
\textbf{\te{పాఠకుడు}} — \emph{pāṭhakuḍu} — \textquotedblleft{}a reader\textquotedblright{}.

\begin{quote}
\te{చాలా} \te{మంది} \te{పాఠకులు} \te{ఉత్తరాలు} \te{రాశారు}.
\end{quote}

\begin{quote}
\emph{cālā mandi pāṭhakulu uttarālu rāśāru.}
\end{quote}

\begin{quote}
Many readers wrote letters.
\end{quote}

Plural: \te{పాఠకులు}, \emph{pāṭhakulu}.

Built on the root of \te{పాఠం}, a lesson: \te{పాఠకుడు}, one who reads.

\begin{cousinweb}
\textbf{\te{పాఠకుడు}} (\emph{pāṭhakuḍu}) is borrowed from Sanskrit \emph{pāṭhaka}, \textquotedblleft{}a reader, a teacher\textquotedblright{} (CDIAL 8036), from \emph{paṭh}, \textquotedblleft{}to read aloud\textquotedblright{}, which also gives \te{పాఠం}, a lesson.
\end{cousinweb}

\begin{grammarlens}[title={What you've built}]
One more word for telling what happened.
\end{grammarlens}

\subsection*{Your turn}
\begin{itemize}
\raggedright
  \item \emph{Say it:} \emph{pāṭhakuḍu}
  \item \emph{Say it:} \emph{pāṭhakuḍu}, once more
  \item \emph{Say it:} \emph{cālā mandi pāṭhakulu uttarālu rāśāru.} — Many readers wrote letters.
  \item \emph{Recall:} say the Telugu for an editor, then the Telugu for an editorial, then say \emph{pāṭhakuḍu} again
\end{itemize}

\begin{tcolorbox}[breakable,colback=teal!4,colframe=teal!35!black,title={Before you move on}]
What does \te{పాఠకుడు} mean? (\textquotedblleft{}A reader\textquotedblright{}.) Say it once more.
\end{tcolorbox}
\section[\texorpdfstring{\te{ధ్రువీకరించు} (dhruvīkariñcu)}{ధ్రువీకరించు (dhruvīkariñcu)}]{\te{ధ్రువీకరించు} (dhruvīkariñcu) — to certify; to confirm}

\label{lesson:TE-C341-dhruvikarincu}



\begin{quote}
Before the new one: say the Telugu for an editorial, then the Telugu for a reader.
\end{quote}

\subsection*{You'll want to know: \te{ధ్రువీకరించు}}
\textbf{\te{ధ్రువీకరించు}} — \emph{dhruvīkariñcu} — \textquotedblleft{}to certify; to confirm\textquotedblright{}.

\begin{quote}
\te{పోలీసులు} \te{ఆ} \te{వార్తను} \te{ధ్రువీకరించారు}.
\end{quote}

\begin{quote}
\emph{pōlīsulu ā vārtanu dhruvīkariñcāru.}
\end{quote}

\begin{quote}
The police confirmed the news.
\end{quote}

Past (I …): \te{ధ్రువీకరించాను}, \emph{dhruvīkariñcānu}.

Built on \te{ధ్రువ}, fixed or firm: to make a report firm. Print often spells it \te{ధృవీకరించు}.

\begin{cousinweb}
\textbf{\te{ధ్రువీకరించు}} (\emph{dhruvīkariñcu}) is built on Sanskrit \emph{dhruva}, \textquotedblleft{}fixed, firm\textquotedblright{} (CDIAL 6892), with \emph{-īkariñcu}, \textquotedblleft{}to make into\textquotedblright{}.
\end{cousinweb}

\begin{grammarlens}[title={What you've built}]
One more word for telling what happened.
\end{grammarlens}

\subsection*{Your turn}
\begin{itemize}
\raggedright
  \item \emph{Say it:} \emph{dhruvīkariñcu}
  \item \emph{Say it:} \emph{dhruvīkariñcu}, once more
  \item \emph{Say it:} \emph{pōlīsulu ā vārtanu dhruvīkariñcāru.} — The police confirmed the news.
  \item \emph{Recall:} say the Telugu for an editorial, then the Telugu for a reader, then say \emph{dhruvīkariñcu} again
\end{itemize}

\begin{tcolorbox}[breakable,colback=teal!4,colframe=teal!35!black,title={Before you move on}]
What does \te{ధ్రువీకరించు} mean? (\textquotedblleft{}To certify; to confirm\textquotedblright{}.) Say it once more.
\end{tcolorbox}
\section[(dialogue)]{Words from chapters 334–341 — a review}

\label{lesson:TE-R341-review-334-341}



\begin{quote}
Say the Telugu for a birth and a newborn baby.
\end{quote}

\subsection*{Your turn}
Say these aloud:

\begin{itemize}
\raggedright
  \item a birth, a newborn baby, a naming ceremony, to take birth, to bless
  \item a crown, a throne, a coronation, to enthrone, a reign
  \item a rumour, hearsay, an exaggeration, an uproar, to disclose
  \item to arrest, to get caught, a prisoner, shackles, the accused
  \item gradually, as time went by, previously, the past, in those days
  \item an army, a soldier, a commander, a march, to march against
  \item instantly, unexpectedly, out of the blue, in the end, to wake up
  \item the main points, an editor, an editorial, a reader, to certify
\end{itemize}

\begin{tcolorbox}[breakable,colback=teal!4,colframe=teal!35!black,title={Before you move on}]
A birth? (\textbf{\te{జననం}}, \emph{jananaṁ}.) A newborn baby? (\textbf{\te{పసికందు}}, \emph{pasikandu}.) A naming ceremony? (\textbf{\te{నామకరణం}}, \emph{nāmakaraṇaṁ}.) To take birth? (\textbf{\te{జన్మించు}}, \emph{janmiñcu}.) To bless? (\textbf{\te{దీవించు}}, \emph{dīviñcu}.) A crown? (\textbf{\te{కిరీటం}}, \emph{kirīṭaṁ}.) A throne? (\textbf{\te{సింహాసనం}}, \emph{siṁhāsanaṁ}.) A coronation? (\textbf{\te{పట్టాభిషేకం}}, \emph{paṭṭābhiṣēkaṁ}.) To enthrone? (\textbf{\te{పట్టం} \te{కట్టు}}, \emph{paṭṭaṁ kaṭṭu}.) A reign? (\textbf{\te{పాలన}}, \emph{pālana}.) A rumour? (\textbf{\te{పుకారు}}, \emph{pukāru}.) Hearsay? (\textbf{\te{వదంతి}}, \emph{vadanti}.) An exaggeration? (\textbf{\te{అతిశయోక్తి}}, \emph{atiśayōkti}.) An uproar? (\textbf{\te{కలకలం}}, \emph{kalakalaṁ}.) To disclose? (\textbf{\te{వెల్లడించు}}, \emph{vellaḍiñcu}.) To arrest? (\textbf{\te{అరెస్టు} \te{చేయు}}, \emph{aresṭu cēyu}.) To get caught? (\textbf{\te{పట్టుబడు}}, \emph{paṭṭubaḍu}.) A prisoner? (\textbf{\te{ఖైదీ}}, \emph{khaidī}.) Shackles? (\textbf{\te{సంకెళ్ళు}}, \emph{saṅkeḷḷu}.) The accused? (\textbf{\te{నిందితుడు}}, \emph{nindituḍu}.) Gradually? (\textbf{\te{క్రమంగా}}, \emph{kramaṅgā}.) As time went by? (\textbf{\te{రానురాను}}, \emph{rānurānu}.) Previously? (\textbf{\te{మునుపు}}, \emph{munupu}.) The past? (\textbf{\te{గతం}}, \emph{gataṁ}.) In those days? (\textbf{\te{అప్పట్లో}}, \emph{appaṭlō}.) An army? (\textbf{\te{సైన్యం}}, \emph{sainyaṁ}.) A soldier? (\textbf{\te{సైనికుడు}}, \emph{sainikuḍu}.) A commander? (\textbf{\te{సేనాధిపతి}}, \emph{sēnādhipati}.) A march? (\textbf{\te{కవాతు}}, \emph{kavātu}.) To march against? (\textbf{\te{దండెత్తు}}, \emph{daṇḍettu}.) Instantly? (\textbf{\te{తక్షణం}}, \emph{takṣaṇaṁ}.) Unexpectedly? (\textbf{\te{అనుకోకుండా}}, \emph{anukōkuṇḍā}.) Out of the blue? (\textbf{\te{ఉన్నట్టుండి}}, \emph{unnaṭṭuṇḍi}.) In the end? (\textbf{\te{తీరా}}, \emph{tīrā}.) To wake up? (\textbf{\te{మేల్కొను}}, \emph{mēlkonu}.) The main points? (\textbf{\te{ముఖ్యాంశాలు}}, \emph{mukhyāṁśālu}.) An editor? (\textbf{\te{సంపాదకుడు}}, \emph{sampādakuḍu}.) An editorial? (\textbf{\te{సంపాదకీయం}}, \emph{sampādakīyaṁ}.) A reader? (\textbf{\te{పాఠకుడు}}, \emph{pāṭhakuḍu}.) To certify? (\textbf{\te{ధ్రువీకరించు}}, \emph{dhruvīkariñcu}.)
\end{tcolorbox}
